\scanpage{269}
en appliquant \(R\varphi_{\unsure{0|\infty},\infty}\)
\[
\underbrace{(R^*\varphi_{N_0,\infty})(\tau_{\infty,\deleted{\unsure{i}}N_0}i)(K)}
\simeq(R^*\varphi_{N_0,\infty})(\tau_{\infty,N_0}L^\bullet)
\]
i.e.
\[
i(K)\simeq\underbrace{(R^*\varphi_{N_0,\infty})(\tau_{\infty,N_0}L^\bullet)}_{\unsure{L'^\bullet|L'}}
\]
\deleted{Je dis que l'hom d'augmentation \(L^\bullet\to\)}

Mais c'est là un isom dans \(D(\mathcal M_{N_0})\) seulement -- rien de plus précis que la donnée de départ -- on aimerait un isom dans \(D(C).\) On va supposer \(i\) pl. fidèle et avec un adjoint à droite \(\kappa,\) donc
\[
\kappa\circ i\simeq\operatorname{id}_C
\]
et on va supposer \(\kappa\) \uline{exact}: Appliquons \(\kappa\) à l'isom obtenu on trouve
\[
K^\bullet\simeq\kappa\bigl(R^*\varphi_{N_0,\infty}(\tau_{\infty,N_0}L^\bullet)\bigr)
\inserted{\simeq(R^*(\kappa\varphi_{N_0,\infty}))(\tau_{\infty,N_0}L^\bullet)}
\]
i.e. \(K^\bullet\) se reconstruit dans \(D(C)\) à l'aide de \(\tau_{\infty,N_0}(i(K^\bullet))\) \deleted{\unsure{donc}} \(\tau_{\infty,N_0}(L^\bullet)\) dans \(D(\mathcal M_\infty).\)
\deleted{On va simplifier l'expression obtenue, en considérant
\[
\varphi_{N_0,\infty}(\tau_{\infty,N_0}(L^\bullet))\longrightarrow
R^*\varphi_{N_0,\infty}(\tau_{\infty,N_0}L^\bullet)
\]
\uline{Lemme} Cet hom est un quasi-isomorphisme

\uuline{Corollaire}
\[
K^\bullet\simeq\kappa\varphi_{N_0,\infty}(\tau_{\infty,N_0}L^\bullet)
\simeq\kappa\varphi_{N_0,\infty}(\tau_{\infty,N_0}iK).
\]}
\margin{268}
